\documentclass[11pt]{article}
\usepackage{icmlko}
\begin{document}

\title{어제 닫은 문을 다른 층에서 다시 열었다 --- 사례 마흔 건은 부호가 아니라 무게를 맞춘다}
\author{월드모델 실험실}
\date{노트 587 · 594 · 2026년 8월 5일}
\maketitle

\begin{abstract}
노트 $582$ 가 덱의 ``새 카테고리 진입 비용 $=$ 사례 $40$ 건''을 재현 못
해 \textbf{다섯 곳에서 지웠다.} 그때 잰 것은 \emph{부호} 였다 --- 짝 라벨
$k$ 건에서 전역 부호를 정하고 뒤집는다. $k{=}40$ 이 $k{=}0$ 과 같았다.

노트 $583$ 이 그 사이 \textbf{앱은 축 무게 순서가 학습과 뒤집혀 있다}를
쟀다(\texttt{entry\_friction} $-0.64$ $>$ \texttt{target\_breadth}
$+0.21$, 학습은 반대). 그러면 맞출 것은 부호가 아니라 \textbf{무게}다.

\begin{center}
\begin{tabular}{lrrrrr}
\hline
 & $k{=}0$ & $k{=}20$ & $k{=}40$ & $k{=}80$ & $k{=}160$ \\
\hline
\textbf{비게임 앱} & $+0.5038$ & $+0.5223$ & $\mathbf{+0.5633}$ & $+0.5986$ & $\mathbf{+0.6132}$ \\
KR 만화 & $+0.6824$ & $+0.6586$ & $+0.6837$ & $+0.6972$ & $+0.7030$ \\
CN 만화 & $+0.3237$ & $+0.2079$ & $+0.3068$ & $+0.3186$ & $+0.2926$ \\
\hline
앱 · 라벨 섞음 & $+0.5038$ & $+0.0690$ & $\mathbf{-0.1521}$ & $+0.0264$ & $-0.0727$ \\
\hline
\end{tabular}
\end{center}

앱 $k{=}40$ 에서 $\mathbf{+0.0595}$ $[+0.0346,+0.0845]$. $k{=}160$ 이면
$+0.6132$ --- 그 짝 자기 자료 천장 $0.6177$ 의 \textbf{$99\%$} 다.
\textbf{사전 등록 넷이 다 섰다.}
\end{abstract}

\begin{figure}[t]
\centering
\includegraphics[width=\textwidth]{figs/f1.pdf}
\caption{\textbf{가} --- 앱만 크게 오르고 라벨을 섞으면 무너진다.
\textbf{나} --- 채택이 옛 결론을 낡게 만들었다.
\textbf{다} --- 같은 ``사례 $40$ 건''인데 층이 다르다.}
\end{figure}

\section{층이 다르다}

\begin{center}
\begin{tabular}{lll}
\hline
 & 무엇을 맞추나 & 결과 \\
\hline
노트 $582$ & $k$ 건에서 \textbf{전역 부호} & $k{=}40 = k{=}0$. \textbf{닫음} \\
\textbf{노트 $587$} & $k$ 건에서 \textbf{축 무게}(능형) & 앱 $\mathbf{+0.0595}$. \textbf{연다} \\
\hline
\end{tabular}
\end{center}

노트 $443$ 의 \texttt{calibrate\_ooc} 도, 노트 $463$ 이 반증한 것도
\emph{부호} 였다. \textbf{무게는 아무도 안 재 봤다.}

\paragraph{대조가 값을 한다.}
같은 $k$ 건에 \textbf{라벨을 섞어} 주면 KR $+0.6824 \to -0.4501$, 앱
$+0.5038 \to -0.6560$. 이득이 절차의 부산물이 아니다.

\paragraph{$k{=}20$ 에서 KR 이 내려간다.}
$+0.6824 \to +0.6586$. \textbf{표본이 적으면 무게 고르기가 과적합한다} ---
노트 $584$ 가 전용 모형에서 본 것과 같은 무늬다(넓은 격자가 교차검증은
올리고 유보는 내린다). $k{=}40$ 에서 회복한다.

\section{그런데 KR 은 이득이 없다}

$+0.0013$ --- $0$ 을 문다. 노트 $583$ 이 답을 미리 냈다: \textbf{KR 은
학습과 무게 순서가 거의 같다.} 고칠 것이 없으면 사례를 줘도 안 고쳐진다.

\begin{quote}
\textbf{약한 시장일수록 사례가 크게 돕는다.} 앱은 $+0.0595$, KR 은
$+0.0013$. 사례 $k$ 건의 값은 \emph{그 시장이 학습과 얼마나 다른가}에
비례하지 \emph{$k$} 에 비례하지 않는다.
\end{quote}

\paragraph{CN 은 내려간다 --- 교란이다.}
$k{=}160$ 이면 $352$ 행 중 \textbf{$45\%$가 채점에서 빠진다.} 사전 등록에
미리 적은 한계다 --- \textbf{$k{=}0$ 대비 차만 읽고 $k$ 끼리는 안 견준다.}

\section{채택이 내 옛 결론을 낡게 만들었다 (노트 594)}

노트 $586$ 채택(\texttt{fixaxes.orient} $+$ \texttt{OOC\_DROP} 해제)으로
집 밖 성적이 올랐으니, 노트 $577 \cdot 579$ 의 천장 표를 다시 재야 한다.
\textbf{분모도 바뀐다} --- 챔피언이 이제 축 다섯을 다 보므로 ``축 전부''
천장이 맞는 분모다.

\begin{center}
\begin{tabular}{lrrr}
\hline
전이 $\div$ 천장 & 노트 $577$ & \textbf{노트 $594$} & \\
\hline
KR 만화 전체 & $0.89$ & $\mathbf{0.95}$ & \\
\textbf{KR 만화 위 절반} & $0.78$ & $\mathbf{1.12}$ & \textbf{천장을 넘는다} \\
비게임 앱 전체 & $0.47$ & $\mathbf{0.82}$ & \\
\textbf{비게임 앱 아래 절반} & $\mathbf{0.22}$ & $\mathbf{0.74}$ & \\
비게임 앱 위 절반 & $0.72$ & $0.95$ & \\
CN 만화 전체 & $0.85$ & $0.88$ & \\
CN 만화 위 절반 & $0.94$ & $0.96$ & \\
\hline
\end{tabular}
\end{center}

\begin{quote}
논문 $446 \cdot 448$ 의 머리가 반쯤 깎인다 --- ``\textbf{앱만 진짜로
전이가 못 간다}''(아래 절반 $0.22$)가 \textbf{$0.74$} 다. 앱 아래 절반
전이가 $+0.0722 \to \mathbf{+0.2418}$ 로 $3.3$ 배다.
\end{quote}

\textbf{그런데 진단은 맞았다.} 노트 $579$ 가 ``\texttt{OOC\_DROP} 이 앱
아래 절반 천장의 $66\%$ 를 버린다''를 쟀고, 되찾으니 \emph{예상한
자리가} 메워졌다. \textbf{틀린 것은 결론이었고 진단은 옳았다.}

그리고 KR 위 절반이 $\mathbf{1.12}$ --- 그 짝 라벨을 한 줄도 안 보고
자기 자료 $5$겹 교차검증을 \textbf{넘는다.}

\section*{덱을 되살렸다}

노트 $582$ 가 지운 ``사례 $40$ 건''을 근거와 함께 되살렸다 ---
\textbf{재학습 없이 $82\%$, 사례 $40$ 건이면 $91\%$, $160$ 건이면
$99\%$}(비게임 앱). 전이 비도 $71\!\sim\!100\% \to \mathbf{82\!\sim\!95\%}$,
상위 절반 $93\!\sim\!192\% \to \mathbf{95\!\sim\!112\%}$ 로 고쳤다.

\begin{quote}
\textbf{지웠다가 되살린 것이 지우지 않은 것보다 낫다.} 사이에 근거가
생겼고, 그 근거는 \emph{다른 층}에서 왔다.
\end{quote}

\section*{조항}

\begin{quote}
\textbf{``안 된다''는 무엇을 맞췄는지까지 적어야 한다.} 노트 $582$ 의
``사례 $40$ 건은 안 오른다''는 \emph{부호를 맞췄을 때} 참이었다.
같은 문장이 무게 층에서는 거짓이다. \textbf{닫은 갈래를 다시 열려면
층을 바꾸면 된다.}
\end{quote}

\section*{남은 것}

\begin{itemize}
\item \textbf{무게가 부호보다 나은가} --- 직접 견주지 않았다. 노트 $443$
\texttt{calibrate\_ooc} 를 같은 $k$ 로 돌려야 안다.
\item \textbf{KR 이 왜 이득이 없나} --- ``이미 무게가 같다''는 읽기는
그럴듯하나 안 쟀다. 무게 순서의 거리를 재서 이득과 붙는지 보면 된다.
\item \textbf{배포에 어떻게 넣나} --- 지금은 실험 대본이다. 짝마다 $k$ 건을
받아 능형을 태우는 자리가 \texttt{lab/} 에 없다.
\end{itemize}

\end{document}
